% Môn: Vật lý
% Lớp: 12
% Chương: Chương I. Vật lí nhiệt
% Tên đề: Ôn tập Chương 1 - Lý 12 - Đề ôn 1 - Tân Bình
% Loại: Trắc nghiệm nhiều lựa chọn + Đúng/Sai + Trả lời ngắn
% Trường: THPT Tân Bình

\section{ÔN TẬP CHƯƠNG I. VẬT LÍ NHIỆT}
\subsection{Đề ôn 1 -- Tân Bình}

\section*{PHẦN I: CÂU TRẮC NGHIỆM PHƯƠNG ÁN NHIỀU LỰA CHỌN}


\right 
\begin{ex}
	% ID: L12C1B1-01-TN
	% Mức: NB
	Chọn câu \textbf{sai} khi nói về cấu tạo chất.
	\choice
	{Các phân tử luôn luôn chuyển động không ngừng.}
	{Các phân tử chuyển động càng nhanh thì nhiệt độ của vật càng cao và ngược lại.}
	{\True Các phân tử luôn luôn đứng yên và chỉ chuyển động khi nhiệt độ của vật càng cao.}
	{Các chất được cấu tạo từ các hạt riêng biệt là phân tử.}
	\loigiai{Các phân tử cấu tạo nên vật luôn chuyển động không ngừng.}
\end{ex}

\begin{ex}
	% ID: L12C1B1-02-TN
	% Mức: NB
	Chỉ ra kết luận \textbf{sai} trong các kết luận sau.
	\choice
	{Sự chuyển từ thể rắn sang thể lỏng gọi là sự nóng chảy.}
	{Sự chuyển từ thể lỏng sang thể rắn gọi là sự đông đặc.}
	{Trong thời gian nóng chảy (hay đông đặc) nhiệt độ của hầu hết các vật không thay đổi.}
	{\True Các kim loại có nhiệt độ nóng chảy giống nhau.}
	\loigiai{Mỗi chất có nhiệt độ nóng chảy riêng nên các kim loại khác nhau có nhiệt độ nóng chảy khác nhau.}
\end{ex}

\begin{ex}
	% ID: L12C1B1-03-TN
	% Mức: NB
	Điều nào sau đây đúng khi nói về cấu tạo chất của thể khí?
	\choice
	{Khoảng cách giữa các phân tử gần nhau.}
	{Sự sắp xếp của các phân tử có trật tự.}
	{\True Các phân tử chuyển động hỗn loạn.}
	{Các phân tử chỉ dao động quanh vị trí cân bằng cố định.}
	\loigiai{Ở thể khí, các phân tử ở cách xa nhau và chuyển động hỗn loạn không ngừng.}
\end{ex}


\dangbt{Đề ôn 2}
\begin{ex}
	% ID: L12C1B1-04-TN
	% Mức: VD
	Một quả bóng khối lượng $100\,\mathrm{g}$ rơi từ độ cao $10\,\mathrm{m}$ xuống sân và nảy lên được $7\,\mathrm{m}$. Lấy $g=9{,}8\,\mathrm{m/s^2}$. Độ biến thiên nội năng của quả bóng trong quá trình trên bằng
	\choice
	{\True $2{,}94\,\mathrm{J}$}
	{$3{,}00\,\mathrm{J}$}
	{$294\,\mathrm{J}$}
	{$6{,}86\,\mathrm{J}$}
	\loigiai{$\Delta U=mg(h_1-h_2)=0{,}1\cdot9{,}8\cdot(10-7)=2{,}94\,\mathrm{J}$.}
\end{ex}


\dangbt{Đề ôn 2}
\begin{ex}
	% ID: L12C1B1-05-TN
	% Mức: NB
	Nội dung nguyên lí I nhiệt động lực học là:
	\choice
	{\True Độ biến thiên nội năng của vật bằng tổng công và nhiệt lượng mà vật nhận được.}
	{Độ biến thiên nội năng của vật bằng tổng nhiệt lượng mà vật nhận được.}
	{Độ biến thiên nội năng của vật bằng tổng công mà vật nhận được.}
	{Độ biến thiên nội năng của vật bằng hiệu số công và nhiệt lượng mà vật nhận được.}
	\loigiai{Theo nguyên lí I nhiệt động lực học: $\Delta U=A+Q$.}
\end{ex}

\begin{ex}
	% ID: L12C1B1-06-TN
	% Mức: TH
	Trường hợp nội năng của vật bị biến đổi \textbf{không phải do truyền nhiệt} là:
	\choice
	{Chậu nước để ngoài nắng một lúc nóng lên.}
	{Gió mùa đông bắc tràn về làm cho không khí lạnh đi.}
	{\True Khi trời lạnh, ta xoa hai bàn tay vào nhau cho ấm lên.}
	{Cho cơm nóng vào bát thì bưng bát cũng thấy nóng.}
	\loigiai{Xoa hai bàn tay vào nhau làm nội năng tăng do thực hiện công.}
\end{ex}

\begin{ex}
	% ID: L12C1B1-07-TN
	% Mức: TH
	Hệ thức $\Delta U=A+Q$ khi $Q<0$ và $A>0$ mô tả quá trình
	\choice
	{hệ truyền nhiệt và sinh công.}
	{hệ nhận nhiệt và sinh công.}
	{\True hệ truyền nhiệt và nhận công.}
	{hệ nhận nhiệt và nhận công.}
	\loigiai{$Q<0$ cho biết hệ truyền nhiệt; $A>0$ cho biết hệ nhận công.}
\end{ex}

\begin{ex}
	% ID: L12C1B1-08-TN
	% Mức: NB
	Mỗi độ chia ($1^\circ\mathrm{C}$) trong thang Celsius bằng $X$ của khoảng cách giữa nhiệt độ tan chảy của nước tinh khiết đóng băng và nhiệt độ sôi của nước tinh khiết ở áp suất tiêu chuẩn. $X$ là
	\choice
	{$\dfrac{1}{273{,}16}$}
	{\True $\dfrac{1}{100}$}
	{$\dfrac{180}{273{,}15}$}
	{$\dfrac{1}{273{,}15}$}
	\loigiai{Khoảng từ $0^\circ\mathrm{C}$ đến $100^\circ\mathrm{C}$ được chia thành $100$ phần bằng nhau nên $X=\dfrac{1}{100}$.}
\end{ex}

\begin{ex}
	% ID: L12C1B1-09-TN
	% Mức: NB
	Điểm đóng băng và sôi của nước theo thang Fahrenheit ở áp suất tiêu chuẩn là
	\choice
	{$0^\circ\mathrm{F}$ và $100^\circ\mathrm{F}$}
	{$100^\circ\mathrm{F}$ và $200^\circ\mathrm{F}$}
	{\True $32^\circ\mathrm{F}$ và $212^\circ\mathrm{F}$}
	{$20^\circ\mathrm{F}$ và $220^\circ\mathrm{F}$}
	\loigiai{$0^\circ\mathrm{C}=32^\circ\mathrm{F}$ và $100^\circ\mathrm{C}=212^\circ\mathrm{F}$.}
\end{ex}


\dangbt{Đề ôn 2}
\begin{ex}
	% ID: L12C1B1-10-TN
	% Mức: TH
	Nhiệt độ trung bình ngày--đêm tại Hà Nội là $24^\circ\mathrm{C}$--$17^\circ\mathrm{C}$. Sự chênh lệch nhiệt độ này trong thang Kelvin là bao nhiêu?
	\choice
	{\True $7\,\mathrm{K}$}
	{$24\,\mathrm{K}$}
	{$17\,\mathrm{K}$}
	{$280\,\mathrm{K}$}
	\loigiai{$\Delta T=24-17=7\,\mathrm{K}$.}
\end{ex}

\begin{ex}
	% ID: L12C1B1-11-TN
	% Mức: TH
	Trong thang nhiệt độ Kelvin, nhiệt độ của nước đá đang tan là $273\,\mathrm{K}$. Hỏi nhiệt độ của nước đang sôi là bao nhiêu?
	\choice
	{$0\,\mathrm{K}$}
	{\True $373\,\mathrm{K}$}
	{$173\,\mathrm{K}$}
	{$100\,\mathrm{K}$}
	\loigiai{$T_{\mathrm{sôi}}=273+100=373\,\mathrm{K}$.}
\end{ex}


\dangbt{Đề ôn 2}
\begin{ex}
	% ID: L12C1B1-12-TN
	% Mức: NB
	Công thức tính nhiệt lượng cần truyền cho vật có khối lượng $m$ để làm vật nóng lên thêm $\Delta T$ độ K là
	\choice
	{\True $Q=mc\Delta T$}
	{$Q=m\Delta T$}
	{$Q=c\Delta T$}
	{$Q=mcT$}
	\loigiai{Nhiệt lượng cần truyền cho vật là $Q=mc\Delta T$.}
\end{ex}

\begin{ex}
	% ID: L12C1B1-13-TN
	% Mức: NB
	Nhiệt hóa hơi riêng là
	\choice
	{lượng nhiệt cần thiết để chuyển $1\,\mathrm{kg}$ chất từ thể rắn sang thể lỏng.}
	{\True lượng nhiệt cần thiết để chuyển $1\,\mathrm{kg}$ chất từ thể lỏng sang thể khí.}
	{lượng nhiệt cần thiết để chuyển $1\,\mathrm{kg}$ chất từ thể khí sang thể lỏng.}
	{lượng nhiệt cần thiết để chuyển $1\,\mathrm{kg}$ chất từ thể lỏng sang thể rắn.}
	\loigiai{Nhiệt hóa hơi riêng là nhiệt lượng cần cung cấp để hóa hơi hoàn toàn $1\,\mathrm{kg}$ chất lỏng ở nhiệt độ sôi.}
\end{ex}

\begin{ex}
	% ID: L12C1B1-14-TN
	% Mức: TH
	Quá trình biến đổi trạng thái của một khối nước đá từ trạng thái rắn sang trạng thái lỏng được biểu diễn bằng đồ thị nhiệt độ theo thời gian. Đoạn đồ thị ứng với quá trình nóng chảy của khối nước đá là
	\choice
	{đoạn $AB$.}
	{\True đoạn $BC$.}
	{đoạn $CD$.}
	{đoạn $AB$ và $CD$.}
	\loigiai{Trong quá trình nóng chảy của nước đá tinh khiết, nhiệt độ không đổi nên đoạn nằm ngang $BC$ biểu diễn quá trình nóng chảy.}
\end{ex}

\begin{ex}
	% ID: L12C1B1-15-TN
	% Mức: TH
	Trong thí nghiệm đo nhiệt hóa hơi riêng của nước, nhiệt hóa hơi riêng của nước được xác định bằng công thức
	\[
	L=\frac{\mathcal{P}\tau}{\Delta m}.
	\]
	Giá trị $\Delta m$ là khối lượng
	\choice
	{nước ban đầu trong ấm đun.}
	{nước trong ấm đun tại thời điểm $\tau$.}
	{\True nước bị bay hơi sau thời gian $\tau$.}
	{nước và ấm đun.}
	\loigiai{Ta có $\mathcal{P}\tau=L\Delta m$, nên $\Delta m$ là khối lượng nước bị bay hơi sau thời gian $\tau$.}
\end{ex}


\dangbt{Đề ôn 2}
\begin{ex}
	% ID: L12C1B1-16-TN
	% Mức: TH
	Biết nhiệt nóng chảy riêng của thiếc là $\lambda=0{,}61\times10^5\,\mathrm{J/kg}$. Nhiệt lượng cần cung cấp cho một cuộn thiếc hàn có khối lượng $0{,}4\,\mathrm{kg}$ nóng chảy hoàn toàn ở nhiệt độ nóng chảy là
	\choice
	{\True $24400\,\mathrm{J}$}
	{$4880\,\mathrm{J}$}
	{$4{,}88\times10^7\,\mathrm{J}$}
	{$76250\,\mathrm{J}$}
	\loigiai{$Q=\lambda m=0{,}61\times10^5\cdot0{,}4=24400\,\mathrm{J}$.}
\end{ex}

\begin{ex}
	% ID: L12C1B1-17-TN
	% Mức: VD
	Tính nhiệt lượng cần cung cấp cho $4\,\mathrm{kg}$ nước đá ở $0^\circ\mathrm{C}$ để chuyển nó thành nước ở $20^\circ\mathrm{C}$. Biết $\lambda=3{,}4\times10^5\,\mathrm{J/kg}$ và $c=4180\,\mathrm{J/(kg\cdot K)}$.
	\choice
	{\True $1694400\,\mathrm{J}$}
	{$136\times10^4\,\mathrm{J}$}
	{$334400\,\mathrm{J}$}
	{$1894400\,\mathrm{J}$}
	\loigiai{$Q=m\lambda+mc\Delta T=4\cdot3{,}4\times10^5+4\cdot4180\cdot20=1694400\,\mathrm{J}$.}
\end{ex}

\begin{ex}
	% ID: L12C1B1-18-TN
	% Mức: VDC
	Người ta dùng một lò nung điện có công suất $20\,\mathrm{kW}$ để làm nóng chảy hoàn toàn $2\,\mathrm{kg}$ đồng có nhiệt độ ban đầu $30^\circ\mathrm{C}$. Biết chỉ $50\%$ năng lượng tiêu thụ của lò được dùng vào việc làm đồng nóng lên và nóng chảy hoàn toàn ở nhiệt độ $1083^\circ\mathrm{C}$. Cho $c=380\,\mathrm{J/(kg\cdot K)}$, $\lambda=1{,}8\times10^5\,\mathrm{J/kg}$. Thời gian cần thiết khoảng
	\choice
	{$58\,\mathrm{s}$}
	{\True $116\,\mathrm{s}$}
	{$90\,\mathrm{s}$}
	{$30\,\mathrm{s}$}
	\loigiai{$Q=mc(1083-30)+m\lambda=1{,}160{,}280\,\mathrm{J}$. Do hiệu suất $50\%$, $Pt=Q/0{,}5$, suy ra $t\approx116\,\mathrm{s}$.}
\end{ex}

\section*{PHẦN II: CÂU TRẮC NGHIỆM ĐÚNG/SAI}


\dangbt{Đề ôn 2}
\begin{ex}
	% ID: L12C1B1-19-DS
	% Mức: NB
	Trả lời đúng/sai khi nói về cấu tạo chất.
	\choiceTF
	{\True Các chất được cấu tạo từ các hạt riêng biệt gọi là phân tử.}
	{Các nguyên tử, phân tử đứng sát nhau và giữa chúng không có khoảng cách.}
	{\True Lực tương tác giữa các phân tử ở thể rắn lớn hơn lực tương tác giữa các phân tử ở thể lỏng và thể khí.}
	{\True Các nguyên tử, phân tử chất lỏng dao động xung quanh các vị trí cân bằng không cố định.}
	\loigiai{Các chất được cấu tạo từ các hạt riêng biệt và giữa chúng có khoảng cách. Lực tương tác ở thể rắn lớn hơn so với thể lỏng và thể khí. Trong chất lỏng, các phân tử dao động quanh các vị trí cân bằng không cố định.}
\end{ex}


\dangbt{Đề ôn 2}
\begin{ex}
	% ID: L12C1B1-20-DS
	% Mức: TH
	Người ta truyền cho khối khí trong xi-lanh một nhiệt lượng $140\,\mathrm{J}$ làm pit-tông đi lên. Khí nở ra và thực hiện công $200\,\mathrm{J}$. Trả lời đúng/sai với các nhận định sau:
	\choiceTF
	{\True Khối khí nhận nhiệt lượng $140\,\mathrm{J}$.}
	{\True Khối khí thực hiện công nên $A<0$ và có giá trị $-200\,\mathrm{J}$.}
	{\True Biểu thức nguyên lí I nhiệt động lực học trong trường hợp này là $\Delta U=A+Q$.}
	{Độ biến thiên nội năng của khí có giá trị là $-340\,\mathrm{J}$.}
	\loigiai{$Q=140\,\mathrm{J}$, $A=-200\,\mathrm{J}$ nên $\Delta U=A+Q=-60\,\mathrm{J}$.}
\end{ex}

\begin{ex}
	% ID: L12C1B1-21-DS
	% Mức: TH
	Nhiệt nóng chảy riêng của chì là $0{,}25\times10^5\,\mathrm{J/kg}$, nhiệt độ nóng chảy của chì là $327^\circ\mathrm{C}$, biết nhiệt dung riêng của chì là $126\,\mathrm{J/(kg\cdot K)}$.
	\choiceTF
	{Miếng chì khối lượng $1\,\mathrm{kg}$ đang ở nhiệt độ $25^\circ\mathrm{C}$ được cung cấp nhiệt lượng $0{,}25\times10^5\,\mathrm{J}$ thì nhiệt độ của nó tăng lên $260^\circ\mathrm{C}$.}
	{\True Cần cung cấp nhiệt lượng $0{,}25\times10^5\,\mathrm{J}$ để làm nóng chảy hoàn toàn $1\,\mathrm{kg}$ chì ở nhiệt độ nóng chảy của nó.}
	{Cần cung cấp nhiệt lượng $25126\,\mathrm{J}$ để làm nóng chảy hoàn toàn $1\,\mathrm{kg}$ chì từ nhiệt độ $25^\circ\mathrm{C}$.}
	{\True Biết công suất của lò nung là $1000\,\mathrm{W}$, giả sử hiệu suất của lò là $100\%$. Thời gian để làm nóng chảy hoàn toàn $1\,\mathrm{kg}$ chì từ nhiệt độ nóng chảy của nó bằng $25\,\mathrm{s}$.}
	\loigiai{a) Sai vì $\Delta T=25000/126\approx198{,}4^\circ\mathrm{C}$ nên nhiệt độ khoảng $223{,}4^\circ\mathrm{C}$. b) Đúng vì $Q=m\lambda=25000\,\mathrm{J}$. c) Sai vì còn phải nung từ $25^\circ\mathrm{C}$ đến $327^\circ\mathrm{C}$: $Q=126(327-25)+25000=63052\,\mathrm{J}$. d) Đúng vì $t=25000/1000=25\,\mathrm{s}$.}
\end{ex}

\begin{ex}
	% ID: L12C1B1-22-DS
	% Mức: VD
	Một học sinh làm thí nghiệm đun nóng để làm $0{,}2\,\mathrm{kg}$ nước đá ở $0^\circ\mathrm{C}$ chuyển hoàn toàn thành hơi nước ở $100^\circ\mathrm{C}$. Cho $\lambda=3{,}34\times10^5\,\mathrm{J/kg}$, $c=4{,}20\,\mathrm{kJ/(kg\cdot K)}$, $L=2{,}26\times10^6\,\mathrm{J/kg}$. Bỏ qua hao phí nhiệt ra môi trường. Trong các phát biểu sau, phát biểu nào đúng, phát biểu nào sai?
	\choiceTF
	{\True Nhiệt lượng cần thiết để làm nóng chảy hoàn toàn $0{,}2\,\mathrm{kg}$ nước đá tại nhiệt độ nóng chảy là $66800\,\mathrm{J}$.}
	{Nhiệt lượng cần thiết để đưa $0{,}2\,\mathrm{kg}$ nước từ $0^\circ\mathrm{C}$ đến $100^\circ\mathrm{C}$ là $8400\,\mathrm{J}$.}
	{Nhiệt lượng cần thiết để làm hóa hơi hoàn toàn $0{,}2\,\mathrm{kg}$ nước ở $100^\circ\mathrm{C}$ là $42500\,\mathrm{J}$.}
	{Nhiệt lượng để làm $0{,}2\,\mathrm{kg}$ nước đá ở $0^\circ\mathrm{C}$ chuyển hoàn toàn thành hơi nước ở $100^\circ\mathrm{C}$ là $60280\,\mathrm{J}$.}
	\loigiai{$Q_1=m\lambda=66800\,\mathrm{J}$; $Q_2=mc\Delta T=84000\,\mathrm{J}$; $Q_3=mL=452000\,\mathrm{J}$. Tổng $Q=602800\,\mathrm{J}$.}
\end{ex}

\section*{PHẦN III: CÂU TRẢ LỜI NGẮN}

\begin{ex}
	% ID: L12C1B1-23-TLN
	% Mức: TH
	Nhiệt độ cơ thể con người là $37^\circ\mathrm{C}$. Xác định nhiệt độ của cơ thể con người trong thang nhiệt độ Kelvin.
	\shortans{310 K}
	\loigiai{$T=37+273=310\,\mathrm{K}$.}
\end{ex}

\begin{ex}
	% ID: L12C1B1-24-TLN
	% Mức: TH
	Cần cung cấp bao nhiêu kJ nhiệt lượng để hóa hơi hoàn toàn $2{,}5\,\mathrm{kg}$ nước đang sôi ở nhiệt độ $100^\circ\mathrm{C}$. Biết nhiệt hóa hơi riêng của nước ở $100^\circ\mathrm{C}$ là $2{,}26\times10^6\,\mathrm{J/kg}$.
	\shortans{5650 kJ}
	\loigiai{$Q=mL=2{,}5\cdot2{,}26\times10^6=5{,}65\times10^6\,\mathrm{J}=5650\,\mathrm{kJ}$.}
\end{ex}

\begin{ex}
	% ID: L12C1B1-25-TLN
	% Mức: VD
	Cung cấp nhiệt lượng $1{,}5\,\mathrm{J}$ cho một khối khí trong một xilanh đặt nằm ngang. Chất khí nở ra, đẩy pit-tông đi một đoạn $5\,\mathrm{cm}$. Độ biến thiên nội năng của khối khí là bao nhiêu Jun? Biết lực ma sát giữa pit-tông và xilanh có độ lớn là $20\,\mathrm{N}$, coi pit-tông chuyển động thẳng đều.
	\shortans{0,5 J}
	\loigiai{Khí thực hiện công chống lực ma sát: $A=-Fs=-20\cdot0{,}05=-1\,\mathrm{J}$. Do đó $\Delta U=Q+A=1{,}5-1=0{,}5\,\mathrm{J}$.}
\end{ex}

\begin{ex}
	% ID: L12C1B1-26-TLN
	% Mức: VDC
	Dùng một bếp điện để đun nóng một nồi đựng $2\,\mathrm{kg}$ nước đá ở $0^\circ\mathrm{C}$. Sau $10$ phút thì nước đá tan chảy hoàn toàn. Nhiệt dung riêng của nước là $4200\,\mathrm{J/(kg\cdot K)}$; nhiệt nóng chảy của nước đá là $3{,}4\times10^5\,\mathrm{J/kg}$. Sau bao nhiêu phút thì nước đá bắt đầu sôi? Kết quả lấy đến $1$ chữ số sau dấu phẩy thập phân.
	\shortans{22,4 phút}
	\loigiai{Trong $10$ phút đầu, bếp cung cấp $Q_1=m\lambda$. Để nước từ $0^\circ\mathrm{C}$ đến $100^\circ\mathrm{C}$ cần $Q_2=mc\Delta T$. Suy ra $t_2=10\dfrac{c\Delta T}{\lambda}=10\dfrac{4200\cdot100}{3{,}4\times10^5}\approx12{,}4$ phút. Tổng thời gian từ lúc đun đến khi nước bắt đầu sôi là $t=10+12{,}4=22{,}4$ phút.}
\end{ex}

\begin{ex}
	% ID: L12C1B1-27-TLN
	% Mức: VD
	Trên một thang đo $X$, điểm băng và điểm sôi lần lượt là $40^\circ X$ và $120^\circ X$. Một thang đo $Y$ khác có điểm băng và điểm sôi tương ứng là $-30^\circ Y$ và $130^\circ Y$. Nếu kết quả đọc được trên thang đo $X$ là $50^\circ X$ thì tương ứng trên thang đo $Y$ là bao nhiêu độ $Y$?
	\shortans{-10 Y}
	\loigiai{Do thang đo tuyến tính: $\dfrac{50-40}{120-40}=\dfrac{Y+30}{130-(-30)}$. Suy ra $Y+30=20$, nên $Y=-10^\circ$.}
\end{ex}

\begin{ex}
	% ID: L12C1B1-28-TLN
	% Mức: VDC
	Một ấm nhôm có khối lượng $0{,}6\,\mathrm{kg}$ chứa $1{,}5\,\mathrm{kg}$ nước ở nhiệt độ $20^\circ\mathrm{C}$, sau đó đun bằng bếp điện. Sau thời gian $35$ phút thì đã có $20\%$ khối lượng nước đã hóa hơi ở nhiệt độ sôi $100^\circ\mathrm{C}$. Biết rằng $75\%$ nhiệt lượng mà bếp cung cấp được dùng vào việc đun nước. Cho $c_{\mathrm{nước}}=4200\,\mathrm{J/(kg\cdot K)}$, $c_{\mathrm{Al}}=880\,\mathrm{J/(kg\cdot K)}$, $L=2{,}26\times10^6\,\mathrm{J/kg}$. Công suất cung cấp nhiệt của bếp là bao nhiêu W? Kết quả đến hàng đơn vị.
	\shortans{777 W}
	\loigiai{$Q=1{,}5\cdot4200\cdot80+0{,}6\cdot880\cdot80+0{,}2\cdot1{,}5\cdot2{,}26\times10^6=1{,}224{,}240\,\mathrm{J}$. Vì $75\%$ nhiệt lượng bếp cung cấp được dùng vào quá trình nên $P=Q/(0{,}75\cdot35\cdot60)\approx777\,\mathrm{W}$.}
\end{ex}
